% =====================================================================================================
\section{The method, as coded, delivers the posterior}\label{sec:delivery}

The previous sections are about the posterior. This section is about the software: what a run actually reports. The main
result, Theorem~\ref{thm:reported-random}, says that there is an event, determined by the data alone and of probability
tending to one, on which every run that produces output reports exactly the true model, whatever its seed. We build up to it in order: the start law already
concentrates on the true model (Section~\ref{sec:start}); the local search then removes what the start mode got wrong
(Section~\ref{sec:reported}). We then show that a certified small model rarely misses a true predictor, under much weaker
conditions (Section~\ref{sec:screen}), and that every reported estimate is consistent (Section~\ref{sec:outputs}).

\subsection{How the signal looks to a block-by-block method}\label{sec:delivery-quantities}

The method inspects the predictors in blocks of at most $16$. The conditions of this section say how the signal looks
from that point of view. They use the following quantities.

\paragraph{Signal gain and residual scale (fixed design).} For a set $U\subseteq\set{p}$ and a predictor $l\notin U$,
the \emph{signal gain} of $l$ given $U$ is $\delta^2_{U,l}:=\norm{(P_{U\cup\{l\}}-P_U)\mu}^2$: the increase in explained
mean signal when $x_l$ joins $X_U$. The \emph{residual scale} left by $U$ is
$\tau^2_U:=\sigs^2+\norm{(I-P_U)\mu}^2/(n-1)$: noise plus unexplained signal energy per observation. The sample
regression coefficient of $x_l$ on $X_U$ is $\hat\beta_{U,l}:=(X_U^\top X_U)^{-1}X_U^\top x_l$, padded with zeros to a
vector indexed by a larger set when needed.

\paragraph{Partial covariances (random design).} With $W\sim\Normal_p(0,\Sig)$, $\Cov(W_j,W_l\mid W_U)$ is the partial
covariance of $W_j$ and $W_l$ given $(W_k)_{k\in U}$; $\bar\beta_{U,l}:=\Sigma_{n,UU}^{-1}\Sigma_{n,Ul}$ is the population
regression coefficient of $W_l$ on $W_U$ (padded with zeros); and $\Var(W_l\mid W_U)$ is the conditional variance.

\paragraph{Visibility and leakage.} A true predictor is \emph{visible} if its signal cannot be explained away by any $15$
other predictors; $15$ is the largest block size minus one. A \emph{leakage set} is a deterministic set $Q$ with
$\Sz\subseteq Q\subseteq\set{p}$ that contains every predictor able to imitate part of the signal. Its \emph{lattice} is
$\Lat(Q):=\{T:\Sz\subseteq T\subseteq Q\}$, the models between the true model and $Q$. Below, \as{NL} is a condition that no null predictor leaks at all, as seen by
the start law; Assumption~\ref{asm:reported-random} allows a leakage set $Q$ larger than $\Sz$, under a different, weaker
leak bound outside $Q$.

\paragraph{Signals fixed on the standardised scale.} Theorems~\ref{thm:start}, \ref{thm:reported-random}
and~\ref{thm:outputs} fix the signal on the scale the method sees:

\begin{assumption}{FSt}{Fixed signals, standardised scale}
The data satisfy $\yr=\alpha\one+X\tstar+\varepsilon$ with $\varepsilon\sim\Normal_n(0,\sigs^2I_n)$, a fixed
$\sigs^2>0$, noise independent of $\Xr$, a deterministic $\tstar$, $\sz\ge1$ fixed, and $\theta_{\mathrm{lo}}\le\abs{\theta^*_j}\le\theta_{\mathrm{hi}}$ for
$j\in\Sz$, with fixed $0<\theta_{\mathrm{lo}}\le\theta_{\mathrm{hi}}<\infty$.
\rulesout{true standardised coefficients that shrink to zero or blow up as $n$ grows.}
\end{assumption}

Since $X=P^\perp_{\one}\Xr D^{-1}$, this is model~\eqref{eq:model} with raw coefficients $\bstar=D^{-1}\tstar$ and an
intercept that depend on the design. In these three theorems \as{FSt} takes the place of \as{GL} and \as{SP}. Raw-scale
versions of the three theorems would need a further argument and are not claimed.

\subsection{The start law already concentrates}\label{sec:start}

The sampler starts from the start law $\pi_0$, which sees interactions within blocks only. When does it already favour
the true model?

\inwords{If every true predictor is visible and no null predictor leaks signal, the start law puts mass $1-n^{-c}$ on the
true model, for every $c<\tfrac12$. If, in addition, null predictors are uncorrelated with the true ones, the start law's
miss is $\OP(\rn)$: as small as the posterior's, up to a factor $\sqrt{\log n}$.}

\begin{assumption}{V}{Visibility}
For every $j\in\Sz$ and every $U\not\ni j$ with $\abs{U}\le15$, $\abs{\Cov(W_j,W^\top\tstar\mid W_U)}\ge a$, for a fixed
$a>0$.
\rulesout{masking: a true predictor whose signal some $15$ other predictors can explain away.}
\end{assumption}

\begin{assumption}{NL}{No leakage, as seen by the start law}
Every null predictor $l\notin\Sz$ has a deterministic \emph{neighbourhood} $\Hl_l\supseteq\{l\}\cup\Sz$ with
$\abs{\Hl_l}\le e^{\kappa_H\sqrt{\log n}}$ for a fixed $\kappa_H$, such that, for every $U\not\ni l$ with $\abs{U}\le15$, writing
$U_0:=U\cap\Hl_l$ for the part of $U$ inside the neighbourhood:
\begin{enumerate}[label=(\roman*),leftmargin=2em,itemsep=1pt]
  \item \emph{no leak:} $\abs{\Cov(W_l,W_j\mid W_{U_0})}\le C_Ln^{-1/2}$ for every $j\in\Sz\setminus U_0$, with a fixed $C_L$;
  \item \emph{locality:} $\abs{\Cov(W_l,W_v\mid W_{U_0})}\le r^\circ_n$ and $\abs{\Cov(W_j,W_v\mid W_{U_0})}\le r^\circ_n$ for every
    $v\in U\setminus\Hl_l$ and $j\in\Sz\setminus U_0$, where $r^\circ_n(n\log n)^{1/4}\to0$.
\end{enumerate}
\rulesout{null predictors that imitate a true one within the start law's block view, and dependence that is not local.}
\end{assumption}

In words: given any block-sized set of predictors, a null predictor is almost uncorrelated with every true predictor, and
only the predictors in its neighbourhood $\Hl_l$ matter for this.

\begin{localassumption}[Small-set floor]\label{asm:start}
$\lmin(\Sigma_{n,S})\ge\varphi_{\mathrm b}^2$ for every $S$ with $\abs{S}\le16+\sz$, for a fixed $\varphi_{\mathrm b}>0$.
\rulesout{population collinearity within sets of the size of a block.}
\end{localassumption}

Assumption~\ref{asm:start} keeps the regressions within a block well posed.

\begin{theorem}[The start law concentrates]\label{thm:start}
Assume \as{GR}, \as{DIM} with $p>\sz$, \as{FSt}, \as{V}, \as{NL} and Assumption~\ref{asm:start}. Then, for every
$c<\tfrac12$,
\begin{equation}\label{eq:start}
  \P\big(\pi_0(\Sz)\ge1-n^{-c}\big)\to1 .
\end{equation}
If moreover every null predictor is uncorrelated with every true predictor, $\Cov(W_l,W_j)=0$ for $l\notin\Sz$ and
$j\in\Sz$, then
\begin{equation}\label{eq:start-rate}
  1-\pi_0(\Sz)=\OP\Big(\sqrt{\frac{\log n}{n}}\Big).
\end{equation}
The exponent $\tfrac12$ in~\eqref{eq:start} cannot be improved: $1-\pi_0(\Sz)\ge\tfrac13(n+1)^{-1/2}$ for every valid data
set with $n\ge4$ and $\sz<K$ (hence for all large $n$). So~\eqref{eq:start-rate} is optimal up to the factor
$\sqrt{\log n}$.
\end{theorem}

\status{\textbf{Proved and independently checked}. Verbatim: Appendix~\ref{app:status}, row 3.3.}

\begin{remark}[Further consequences]\label{rem:start-cons}
\emph{With leakage.} Keep the other assumptions of Theorem~\ref{thm:start}, and let \as{NL} hold with $\Sz$ replaced by a
deterministic $Q\supseteq\Sz$ with $\abs{Q}\le K$, so that only predictors outside $Q$ are required not to leak. Then
$\P\big(\pi_0(\Lat(Q))\ge1-(\abs{Q}+1)n^{-1/2}e^{C\sqrt{\log n}}\big)\to1$; for $\abs{Q}\le n^{1/2-\epsilon}$ this is at
least $1-n^{-c}$ for every $c<\epsilon$.
\end{remark}

\begin{remark}[On the assumptions]\label{rem:start-asm}
(i) \as{V} says that each true predictor, alone, improves the fit of every small model it is missing from; then the start
law, which sees only within-block interactions, still ranks $\Sz$ first.
(ii) In \as{NL} the level $n^{-1/2}$ is a threshold of the proof (numerics show no cliff there); localisation is needed
because the blocks depend on the data. The size bound on $\Hl_l$ enters only through a union bound over these sets;
it is not known to be needed.
(iii) The uncorrelatedness condition for~\eqref{eq:start-rate} is an additional hypothesis for the sharper rate only; it
does not replace \as{NL}. It suffices that these correlations be at most $C/\sqrt{n\log n}$ given every proper subset of
the true model. It is not known to be necessary.
\end{remark}

Theorem~\ref{thm:start} is about an internal object of the method: it explains why the sampler starts in the right
place.

\subsection{The reported model is the true model}\label{sec:reported}

This is the main result. We state it first for Gaussian rows, with conditions on the population law, and then for fixed
designs, with conditions on the sample.

\begin{example}[Visibility and leakage with three predictors]\label{ex:visibility}
Consider three of the $p$ predictors. Let the true model be $\Sz=\{1,2\}$ with $\tstar_{\Sz}=(1,-1)$, let
$\Corr(W_1,W_2)=\varrho$, and let $W_3$ be a null predictor. With $U=\emptyset$,
$\Cov(W_1,W^\top\tstar)=\Cov(W_1,W_1-W_2)=1-\varrho$. So \as{V} requires, among other things, $1-\varrho\ge a$: as
$\varrho\to1$ the two true predictors cancel, and each becomes invisible. Now let $\Corr(W_3,W_1)=\varrho'$ be fixed and
nonzero ($\Corr(W_3,W_2)$ may take any admissible value). Then $W_3$ can imitate part of the signal of $W_1$, so the
leakage conditions below require $3\in Q$: a null predictor that stays correlated with a true one belongs to the leakage
set. Such a design violates \as{NL}, so Theorem~\ref{thm:start} does not apply to it, while
Theorem~\ref{thm:reported-random} can.
\end{example}

\paragraph{Random Gaussian design.}

\inwords{If every true predictor is visible, the predictors that can imitate the signal form a set $Q$ of at most $n^b$
predictors with $b<\tfrac12$, and regressions within $Q$ are local, then, with probability tending to one, every run that
produces output reports the true model. The start mode contains every true predictor and possibly some members of $Q$; the
local search from the start mode then removes the spurious ones, one per move.}

The delivery clauses of the theorem refer to the domain where the code's two fixed integer limits on the local-search move
caps are not reached:
\[
  \Dom:=\big\{(n,\praw):\lceil\log_2(n+\praw+1)\rceil\le10^6,\ \ \lceil\log_2(n+\praw+1)\rceil+\lfloor n/2\rfloor\le2^{31}-1\big\}.
\]
The limits are a ceiling of $10^6$ on the base move cap $\ell_n$, and the 32-bit integer bound $2^{31}-1$ on the move cap
$\ell_n+\abs{\tSz}$ of the search from the start mode. The first can bind only when $n+p+1$ exceeds $2^{10^6}$; the second only when the start mode has more than
about $2.1\times10^9$ predictors, which needs $n\gtrsim4.3\times10^9$, since $\abs{\tSz}\le\lfloor n/2\rfloor$.

Fix $b\in(0,\tfrac12)$ and $b_H>0$ with $15b_H+b<\tfrac12$.

\begin{localassumption}[Polynomial leakage set with local regressions]\label{asm:reported-random}
\begin{asmitems}
  \item\label{asm:rr-leak} \emph{Leakage set.} There is a deterministic $Q\supseteq\Sz$ with $\abs{Q}\le n^b$ and, for
    large $n$, $\abs{Q}\le K$ and ($Q=\set{p}$ or $K\ge2\abs{Q}+32\sz+1$), such that: (i) \emph{no leak outside $Q$:}
    there is a deterministic $\varpi_n$ with $\varpi_n\sqrt{n/\log n}\to0$ such that
    $\abs{\Cov(W_l,W_j\mid W_U)}\le\varpi_n$ for all $l\notin Q$, $j\in\Sz\setminus U$ and $\abs{U}\le15$; and (ii)
    \emph{locality:} there is a deterministic $r_n\to0$, and every $l\notin Q$ has a deterministic
    $\Hl_l\supseteq\{l\}\cup\Sz$ with $\abs{\Hl_l}\le n^{b_H}$, such that for all $l\notin Q$, all $U\not\ni l$ with
    $\abs{U}\le15$ and $U_0:=U\cap\Hl_l$: $\abs{\Cov(W_l,W_v\mid W_{U_0})}\le r_n$ and
    $\abs{\Cov(W_j,W_v\mid W_{U_0})}\le r_n$ for all $v\in U\setminus\Hl_l$ and $j\in\Sz\setminus U_0$.
  \item\label{asm:rr-floor} \emph{Floor on the leakage set.} $\lmin(\Sigma_{n,Q\cup\{l\}})\ge\varphi^2$ for every
    $l\in\set{p}$, with $\varphi$ as in \as{EFS}.
  \item\label{asm:rr-local} \emph{Local regressions within the leakage set.} There are deterministic sequences
    $d_n=o(\log n)$ and $L_n$, and for every $l\in\set{p}$ a deterministic $N_l\subseteq Q\setminus\{l\}$ with
    $\abs{N_l}\le d_n$, such that $\norm{\bar\beta_{T,l}-\bar\beta_{T\cap N_l,l}}_1\le L_n$ for every $T\in\Lat(Q)$ with
    $l\notin T$, where $L_n\to0$ and $L_n^2\log(\abs{Q}+1)=o(\log n)$.
\end{asmitems}
\rulesout{(a) many predictors that can imitate the signal (those that can must form a set of at most $n^b$ predictors,
$b<\tfrac12$); (b) collinearity between the leakage set and any other predictor; (c) regressions, inside the leakage set,
that depend on many of its members at once.}
\end{localassumption}

\begin{theorem}[The reported model equals the true model: random Gaussian design]\label{thm:reported-random}
Let $\Sig$ be positive definite and assume \as{GR}, \as{DIM}, \as{FSt}, \as{V}, \as{EFS}, Assumption~\ref{asm:reported-random}
and Convention~\as{C1}. Then, under the joint law of the design and the noise, there are data events $\mathcal G_n$ with
$\P(\mathcal G_n)\to1$ on which \emph{every} admissible run that produces output reports $\Shat=\Sz$, whatever its seed and the
outcome of its internal calibration. In particular,
\begin{equation}\label{eq:reported}
  \P(\text{the run produces output and }\Shat\neq\Sz)\to0 .
\end{equation}
If moreover Convention~\as{C2} holds, then the run produces output with probability one for every admissible $n\ge3$
with $(n,\praw)\in\Dom$. Under Convention~\as{C2p} instead,
\begin{equation}\label{eq:delivery}
  \P(\text{the run produces output})=1\ \ \text{for every admissible }n\ge3,\qquad \P(\Shat=\Sz)\to1 .
\end{equation}
\end{theorem}

Under \as{C2} the delivery clause holds only on the bounded domain $\Dom$, so no limit in $n$ can be stated there;
\as{C2p} removes the bound. In practice $\Dom$ contains every feasible $(n,\praw)$: see the paragraph before the
theorem.

\status{\textbf{Proved and checked}; its one condition (the transfer from population to sample conditions) is itself
proved and independently checked. Delivery clauses: \textbf{proved and independently checked}. \bridgepending{} A
revision of this statement, from a strengthening that has passed its review, is pending. Verbatim:
Appendix~\ref{app:status}, row 3.1.}

\begin{remark}[Further consequences]\label{rem:reported-random-cons}
(i) On $\mathcal G_n$ the local search from the start mode only deletes the spurious members of the start mode, one per
move, in at most $n^b$ moves, and certifies $\Sz$; its move cap never binds.
(ii) The case $\abs{Q}\le(\log n)^{\kappa'}$ for a fixed $\kappa'$ is a special case.
(iii) Designs that satisfy the assumptions include correlation that is local in a metric, polynomial correlation decay
with exponent greater than $2$ (not yet checked), and one-factor blocks; these are examples, not hypotheses.
\end{remark}

\begin{remark}[On the assumptions]\label{rem:reported-random-asm}
(i) \emph{Why $b<\tfrac12$.} At model size about $\sqrt n$, the prior-plus-Occam penalty of one more predictor
($\approx\log p+\tfrac12\log n$) stops dominating the best of $p$ noise gains ($\approx\log p$). For larger $b$ the key
events of the argument are expected to fail (heuristic; whether the conclusion fails is open).
(ii) Assumption~\ref{asm:reported-random}\ref{asm:rr-local} is a sufficient form of a width condition that the argument needs; it is not known to
be necessary for the conclusion.
(iii) All assumptions are on the population law. The leakage and locality conditions are long and specific; they are the
price of a statement about a \emph{search} algorithm: the signal is not masked, what can imitate it is confined to a set
of size $n^b$, and within that set regressions are local.
\end{remark}

\begin{remark}[Comparison with \citet{SongLiang2015}]\label{rem:SL-reported}
This is consistency of what the software reports, with $p\gg n$ and without assuming a particular correlation
structure for the design. We know of no comparable result for
the selection output of an SMC or MCMC implementation of Bayesian variable selection: \citet{SongLiang2015} run the stochastic
approximation Monte Carlo (SAMC) sampler of \citet{LiangLiuCarroll2007}, citing its general properties, but prove nothing about its selection output.
\end{remark}

\paragraph{Fixed design.} The same conclusion holds for a fixed design under sample-level conditions.

\inwords{If the posterior conditions of Theorem~\ref{thm:post-fixed} hold with strong signals, every true predictor
explains much more than $\log n$ of the signal given any $15$ others, predictors outside a leakage set $Q$ never look like
signal, and $Q$ is small or its regressions are local, then every run that produces output reports the true model, with
probability tending to one.}

\begin{localassumption}[Fixed-design visibility and leakage]\label{asm:reported-fixed}
\begin{asmitems}
  \item\label{asm:rf-floor} \emph{Floors.} $\lmin(\Rhat_T)\ge\varphi_0^2$ for every $T$ with $\abs{T}\le16$ and for every
    $T=Q\cup\{l\}$, $l\in\set{p}$, with $Q$ as in~\ref{asm:rf-leak}.
  \item\label{asm:rf-vis} \emph{Visibility.} There is a deterministic $\omega_n\to\infty$ such that
    $\delta^2_{U,j}\ge\omega_n\log n$ for every $j\in\Sz$ and every $U\not\ni j$ with $\abs{U}\le15$.
  \item\label{asm:rf-leak} \emph{Bounded leakage.} There are a deterministic $Q\supseteq\Sz$ (with $\abs{Q}\le K$ and
    ($Q=\set{p}$ or $K\ge2\abs{Q}+32\sz+1$) for large $n$) and a fixed $\eta\in[0,\eta_A)$,
    $\eta_A:=(\sqrt{2A+1}-\sqrt{2A})^2$, such that $\delta^2_{U,l}\le\eta\,\tau^2_U\log n$ for every $l\notin Q$ and every
    $U\not\ni l$ with $\abs{U}\le15$.
  \item\label{asm:rf-size} \emph{Size of the leakage set.} One of: (d1) $\abs{Q\setminus\Sz}\le(1-\delta_Q)\log n$ for a
    fixed $\delta_Q\in(0,1)$; or (d2) $\abs{Q}\le n^b$ with $b<\tfrac12-\tfrac\eta2-\sqrt{2A\eta}$,
    $\theta^{*2}_{\min}n/(\abs{Q}+\log n)\to\infty$, and local regressions: for every $l$ a deterministic set
    $N_l\subseteq Q\setminus\{l\}$ with $\abs{N_l}\le h\log_2n$, for a fixed $h$ with $h+b<\tfrac12$, and a deterministic $L_n\to0$ with
    $L_n^2\log(\abs{Q}+1)=o(\log n)$, such that $\norm{\hat\beta_{T,l}-\hat\beta_{T\cap N_l,l}}_1\le L_n$ for every
    $T\in\Lat(Q)$ with $\abs{T}<K$ and $l\notin T$.
\end{asmitems}
\rulesout{(a) collinearity within a block, or between the leakage set and any other predictor; (b) masking in the
sample; (c) predictors outside $Q$ whose signal gain, given a small model, exceeds $\eta\,\tau^2_U\log n$; (d) a leakage set so large, or so entangled, that adding a predictor could improve the score
somewhere between $\Sz$ and $Q$.}
\end{localassumption}

\begin{theorem}[The reported model equals the true model: fixed design]\label{thm:reported-fixed}
Let the design be fixed, with no raw predictor constant in the sample (an input condition of the delivery result; then
$p=\praw$). Assume \as{GL}, \as{DIM}, \as{SP}, \as{BC},
\as{EF}, \as{NC}, \as{BM}(b), Assumption~\ref{asm:reported-fixed} and Convention~\as{C1}. Then the conclusions of
Theorem~\ref{thm:reported-random} hold: there are data events $\mathcal G_n$ with $\P(\mathcal G_n)\to1$ on which every
admissible run that produces output reports $\Shat=\Sz$, for every seed and calibration outcome; and the delivery
clauses hold under \as{C2} on $\Dom$ and under \as{C2p} for every admissible $n\ge3$.
\end{theorem}

\status{\textbf{Proved and independently checked}. \bridgepending{} A revision of this statement, from a strengthening
that has passed its review, is pending. Verbatim: Appendix~\ref{app:status}, row 3.2.}

\begin{remark}[On the assumptions]\label{rem:reported-fixed}
(i) \emph{Intuition.} Visibility: each true predictor, added alone to any small model that misses it, explains much more
than $\log n$ of the signal, more than the prior charges for one predictor; so the start law never leaves a true predictor
out. Bounded leakage: outside $Q$ no predictor looks like signal, so none enters the start mode. Size of $Q$: predictors
in $Q$ may enter the start mode, and the local search removes them one at a time; this needs that the score never rewards
\emph{adding} a predictor between $\Sz$ and $Q$, which holds when $Q$ has fewer than $\log n$ extra predictors
(case~(d1)), or when regressions on such models are carried by a few members of $Q$ (case~(d2)).
(ii) These are sample-level conditions on $(X,\tstar)$. They are interpretable but not reducible to eigenvalue floors and
beta-min: visibility fails on well-conditioned designs where true predictors cancel. The proof of
Theorem~\ref{thm:reported-random} shows that population analogues imply them with probability tending to one.
\end{remark}

\begin{remark}[What is not claimed]\label{rem:reported-not}
No result is available without any condition on the lattice of the leakage set (that would need control of the actual
search path), or for any $b\ge\tfrac12$.
\end{remark}

\subsection{Certified screening}\label{sec:screen}

The previous theorems need visibility and bounded leakage. The next result needs neither: it says only that when a run
certifies a small model, that model rarely misses a true predictor. Its conditions can hold on designs, such as designs
with strong factors, where those of Section~\ref{sec:reported} fail.

\inwords{If the design has a sparse eigenvalue floor, the signals and coefficients are bounded by fixed constants, and the
size budget $k_n$ is below the model-size cap and not too large, then the probability that a run produces output,
certifies a reported model of at most $k_n$ predictors, and that model misses a true predictor tends to zero. No
visibility, leakage, start or block condition is used.}

In this subsection $k_n$ is a model-size budget chosen by the analyst.

\begin{localassumption}[Certified screening]\label{asm:screen}
\begin{asmitems}
  \item\label{asm:screen-floor} \emph{Eigenvalue floor.} $\lmin(\Rhat_S)\ge\kappa$ for every $S$ with $\abs{S}\le k_n+\sz$,
    for a fixed $\kappa>0$.
  \item\label{asm:screen-signal} \emph{Signals and coefficients bounded by fixed constants.} $\tmin\ge\tlow>0$ and
    $\norm{\tstar}_1\le M_1$, where $\kappa$, $\tlow$, $M_1$ and $\sigs^2$ do not depend on $n$.
  \item\label{asm:screen-budget} \emph{Size budget.} $k_n\in\{0,\dots,K-1\}$ and $(k_n+1)\log p+\log n\le c_{\mathrm{SB}}\,n$,
    where $c_{\mathrm{SB}}:=\kappa^2\tlow^2/\{200(M_1^2+\sigs^2)\}$.
\end{asmitems}
\rulesout{(a) collinearity among small sets of predictors; (b) signals that vary with $n$ (see
Theorem~\ref{thm:screen-margin}); (c) a budget $k_n$ too large for a union bound over models of that size.}
\end{localassumption}

\begin{theorem}[Certified screening]\label{thm:screen}
Let the design be fixed and assume \as{GL}, Assumption~\ref{asm:screen} and Convention~\as{C1}; \as{SP} is not assumed.
Then
\begin{equation}\label{eq:screen}
  \P\big(\text{the run produces output},\ \Shat\text{ is certified},\ \abs{\Shat}\le k_n,\ \Sz\not\subseteq\Shat\big)\to0 .
\end{equation}
In particular this holds for $p\le n^A$ and every $k_n=o(n/\log n)$.
\end{theorem}

\status{\textbf{Proved and independently checked}. Verbatim: Appendix~\ref{app:status}, row 3.7.}

No condition on the start, leakage, blocks or visibility is used, and $p\le n^A$ is not needed. The proof needs only a
weaker floor, along certain directions, which leaves correlations among null predictors free; the full floor
in Assumption~\ref{asm:screen}\ref{asm:screen-floor} does constrain them.

Can the signal-strength conditions vary with $n$? The next result replaces them by a single margin condition on how much
signal a missing true predictor can add.

\inwords{If every small model that misses a true predictor has some missing true predictor whose gain beats the noise
level of a union bound over models of that size, the conclusion of Theorem~\ref{thm:screen} holds, with no condition on
the dimension, the true model size, or the null predictors.}

For a model $S$ let $\mathcal R_S:=I_n-X_S(G_S+\lambda_0I)^{-1}X_S^\top$, with $\lambda_0:=1/(n+1)$, be the ridge residual
operator ($\mathcal R_\emptyset=I_n$), and for $t\notin S$ let $a(t\mid S):=x_t^\top\mathcal R_Sx_t+\lambda_0$. The
\emph{noiseless gain} of $t$ given $S$ is
\[
  \gamma(t\mid S):=\frac{(x_t^\top\mathcal R_S\mu)^2}{a(t\mid S)}=\frac{n+1}{n}\big(\SSq^{\mu}_S-\SSq^{\mu}_{S\cup\{t\}}\big),
\]
where $\SSq^{\mu}_S$ is the shrunken sum of squares $\SSq_S$ of~\eqref{eq:marglik} with $y$ replaced by its mean $\mu$: the
decrease of the method's own quadratic form when $t$ joins $S$ and the response is replaced by its mean. Let
$\Gamma_n(q):=\min\{\max_{t\in\Sz\setminus S}\gamma(t\mid S):\abs{S}\le q,\ S\not\supseteq\Sz\}$, so that every model of
size at most $q$ that misses a true predictor has some missing true predictor with gain at least $\Gamma_n(q)$; and let
$\tau_n^2:=\sigs^2+\norm{\mu}^2/(n-1)$, which is $\tau^2_\emptyset$ of Section~\ref{sec:delivery-quantities}.

\begin{theorem}[Certified screening with signals that may vary with $n$]\label{thm:screen-margin}
Let the design be fixed and assume \as{GL} and Convention~\as{C1}. If $k_n\in\{0,\dots,K-1\}$ and
\begin{equation}\label{eq:gain-margin}
  \liminf_{n\to\infty}\frac{\sqrt{\Gamma_n(k_n)}}{\tau_n\sqrt{2\log p+\log n}+\sigs\sqrt{2k_n\log p}}>1 ,
\end{equation}
then~\eqref{eq:screen} holds. No condition $p\le n^A$, no bound on $\sz$, and no condition on the start, leakage, blocks
or the collinearity of the null predictors is used.
\end{theorem}

\status{\textbf{Proved and independently checked} (a non-asymptotic form, with probability at most $3e^{-x}$ under a
finite-$n$ margin, is also checked). Verbatim: Appendix~\ref{app:status}, row 3.7.}

\begin{remark}[What is not claimed]\label{rem:screen}
(i) The bound is on the probability of an intersection: a run produces output, certifies a small model, \emph{and} that
model misses a true predictor. It is not a guarantee conditional on certification, and it does not say that the method
returns a superset of the true model with probability tending to one.
(ii) A run is not claimed to certify a small model with probability tending to one; how often it does is reported from
numerics below.
(iii) It is one-sided: certified small models are not shown to equal the true model. Exact recovery from the certificate
needs a further, stronger condition on the design (proved and checked), which we do not state here.
\end{remark}

\paragraph{How often runs certify (numerics, not a theorem).} In existing numerics the method certified the reported model
in every completed run: 2,394 of 2,394 (the completed runs) in the 2,400-fit study ($n=200$); 376 of 376 in the strong-factor experiment
($n\le150$, $p$ up to $10n$); and 71 of 71 in the screening study ($n=100$). Where model sizes were recorded, the certified
models were small: all 376 strong-factor fits had at most $2\sz$ predictors, and all 71 screening fits at most 5 (the
2,400-fit study records no per-fit size). On strong-factor designs (equicorrelation $\varrho\in\{0.5,0.7\}$, $p\ge2.5n$,
$\sz\in\{4,8\}$) the certified model contained every true predictor in 61 of 122 fits at $n=50$, 18 of 20 at $n=75$ and
30 of 30 at $n=100$. At the constants $\kappa$, $\tlow$, $M_1$, $\sigs$ of these designs,
Assumption~\ref{asm:screen}\ref{asm:screen-budget} requires $n\gtrsim1{,}450$. [The three studies are described in the
simulation section of the paper.]

\subsection{The reported estimates are consistent}\label{sec:outputs}

The method reports an evidence estimate, pooled inclusion probabilities and the corresponding median probability model.
This subsection shows that all of them are consistent, on the event that the method's calibration succeeds.

This result concerns runs that use the sampler, that is, runs in which the model space is too large to enumerate; the
enumeration case is not covered. Let $\Acal$ be the event that the method uses the sampler and its calibration succeeds
(possibly through a fallback, Section~\ref{sec:method}); every completed sampler run lies in $\Acal$. On $\Acal$, let $\Zhat$, $\PIPhat$ and $\MPMhat$ be the
evidence estimate, the pooled inclusion probabilities and their median probability model, the values the four islands would produce, pooled with equal weights (Section~\ref{sec:method}), if the run
were not interrupted. On completed runs they are exactly the reported values. Let $\rho_n:=n^{-1/2}e^{C_3\sqrt{\log n}}$ for a constant $C_3$ that depends only on the fixed quantities of the
assumptions; so $\rho_n\le n^{-c}$ eventually, for every $c<\tfrac12$. Assumption~\ref{asm:start} is not listed in the
theorem below because \as{EFS} implies it: $16+\sz\le5(A+1)(\sz+1)\log n$ for $n\ge3$.

\inwords{Under the conditions that make the start law concentrate, every reported estimate is within $\rho_n$ of its
target on the event that the calibration succeeds: the evidence estimate is close to the evidence, the inclusion
probabilities are close to $1$ for true predictors and to $0$ for null ones, and the reported median probability model
and reported model are wrong with vanishing probability.}

\begin{theorem}[The reported estimates are consistent]\label{thm:outputs}
Assume \as{GR}, \as{DIM} with $p>\sz$, \as{FSt}, \as{EFS}, \as{V}, \as{NL}, and Conventions~\as{C1} and~\as{C3}. Then,
under the joint law of the design, the noise and the method's random numbers,
\begin{gather}
  \ind{\Acal}\Big(\big|\Zhat/Z-1\big|+\big|\log\Zhat-\log Z\big|+\max_{j\in\set{p}}\big|\PIPhat_j-\ind{j\in\Sz}\big|\Big)
    =\OP(\rho_n),\label{eq:outputs-rate}\\
  \P\big(\Acal\cap\{\MPMhat\neq\Sz\}\big)\to0,\qquad \P\big(\text{the run completes and }\Shat\neq\Sz\big)\to0,
    \label{eq:outputs-sel}
\end{gather}
and~\eqref{eq:outputs-rate} also holds for the reported values multiplied by the indicator of a completed run.
\end{theorem}

\status{\textbf{Proved and independently checked}. \bridgepending{} Verbatim: Appendix~\ref{app:status}, row 3.6.}

\begin{remark}[What is not claimed]\label{rem:outputs}
(i) No lower bound on $\P(\Acal)$ is asserted: the statements are about intersections with $\Acal$ or with completion, not
conditional on them. Under Convention~\as{C2} completion holds on $\Dom$, and under \as{C2p} for every $n$; that link is
not used here.
(ii) The rate $\rho_n$ is the start law's rate (Remark~\ref{rem:start-cons} with $Q=\Sz$). It is slower than the posterior
rate: $\rho_n/\rn=e^{C_3\sqrt{\log n}}/\sqrt{\log n}\to\infty$.
\end{remark}

\begin{remark}[On the assumptions]\label{rem:outputs-moderate}
Asymptotically the start law concentrates on $\Sz$, so the sampler has almost no mass to move and $\Zhat/Z\to1$, whatever
the length of the sequence of intermediate distributions. At moderate $n$ the posterior is spread over $\Sz$ and its many
one-predictor supersets (Remark~\ref{rem:moderate}). An alternative argument, which does not go through the evidence estimate, gives the same
conclusions for the inclusion probabilities and the median and highest posterior models under a global-overlap
condition (stated in full with the proof).
\end{remark}

\status{The alternative argument: \textbf{proved and independently checked}. Verbatim: Appendix~\ref{app:status}, row 3.6.}

\begin{remark}[The pooled evidence estimate is unbiased]\label{rem:unbiased}
Under Convention~\as{C3}, given the calibration, each island's evidence estimate is unbiased for $Z$, for every design and
data set; no condition
on the design is needed. By linearity, so is the reported equal-weight pooled estimate: on $\Acal$,
$\E[\Zhat\mid\text{calibration}]=Z$. This is the classical unbiasedness of sequential Monte Carlo estimates of normalising
constants \citep{DelMoral2004}, established here for the method as coded, calibration included: the calibration is frozen
before the final run's random numbers are drawn.
\end{remark}

\status{Per island: \textbf{proved and independently checked}. Pooled form: \textbf{proved (a two-line argument), not yet
independently checked}. Verbatim: Appendix~\ref{app:status}, row 3.5.}
